\documentclass[12pt, letterpaper]{article}
\usepackage[utf8]{inputenc}
\usepackage{mathptmx}
\usepackage{geometry}
\usepackage{amsmath, amssymb}
\usepackage{booktabs}
\usepackage[style=apa, backend=biber]{biblatex}
\addbibresource{references.bib}
\geometry{letterpaper, margin=1in}
\usepackage{setspace}
\onehalfspacing
\usepackage{enumitem}
\setlist{itemsep=2pt, parsep=0pt, topsep=4pt}
\usepackage{titlesec}
\usepackage{fancyhdr}
\pagestyle{fancy}
\setlength{\headheight}{15pt}
\fancyhead[L]{Executive Summary}
\fancyhead[R]{\thepage}
\fancyfoot{}
\renewcommand{\headrulewidth}{0pt}

\begin{document}

\begin{center}
    {\Large \textbf{Executive Summary}} \\[0.5em]
    {\large Hybrid Quantum-Classical Software Stack for HPC Algorithm Acceleration} \\[0.3em]
    {Anna Payne} \\[0.3em]
    {IE University -- Capstone Thesis, 2026}
\end{center}

\vspace{0.5em}

\subsection*{Context and Motivation}
The recent integration of Quantum Processing Units (QPUs) within High Performance Computing (HPC) environments
shows great potential for quantum hardware to accelerate previously computationally 
intractable tasks. In the era of Noisy Intermediate-Scale Quantum (NISQ) hardware, hybrid quantum-classical 
algorithms such as the VQE offer a promising path, where QPUs handle circuit executions while 
classical resources manage the heavy optimization loop and preprocessing. However, there is an existing architectural 
gap, as classical HPC frameworks (MPI, CUDA) are not developed for communication with quantum control 
stacks and likewise, quantum frameworks (Qiskit, Cirq) lack distributed HPC integration. Without an available bridging 
middleware, the potential speedup of hybrid systems is negated by several communication bottlenecks.
This thesis presents a hybrid quantum-classical middleware stack targeting four objectives, including  
an accessible API that abstracts the cloud QPU backend, 
MPI distributed parallelization of the Pauli term evaluations, 
GPU-accelerated statevector construction with CUDA and cuStateVec, and lastly 
a performance analysis across the execution paths of CPU, GPU, and QPUs.

\subsection*{Methodology}
This middleware implements a triple-heterogeneous architecture with layers for Orchestration (the MPI manager worker 
topology for parameter distribution and energy aggregation), Acceleration (CUDA kernels with 
mixed precision FP32/FP64 for GPU-accelerated statevector construction), and the Quantum Interface 
(dual path dispatcher, C++ MPI and Python EstimatorV2 paths access the IBM Quantum cloud QPU). 
Stack validation uses the VQE algorithm with a Hardware Efficient (HWE) ansatz and SPSA classical
optimizer on the four benchmark molecules $H_2$, $LiH$, $BeH_2$, $H_2O$ under the STO-3G basis set alongside 
FCI ground truth references. All simulator experiments ran on a single Docker host (Intel i7-1065G7, 
32~GB RAM, NVIDIA GTX 1650 Mobile). QPU experiments used IBM's \texttt{ibm\_marrakesh} 156-qubit Heron processor.

\subsection*{Key Findings}

\textbf{1. MPI distribution provides measurable speedup for complex molecules.} The distributed CPU only 
stack achieved a $1.18\times$ speedup on molecule $H_2O$ (1086 Pauli terms) at $P=2$ ranks compared to the 
serial baseline, however, at $P=8$ regressed to $0.24\times$ due to shared memory contention on a single host. The 
key finding is that on a single host, MPI distribution alone is insufficient without added GPU acceleration.

\textbf{2. GPU acceleration resolves CPU scaling bottleneck.} With GPU Aer, $H_2O$ at $P=8$ 
provided a $1.51\times$ speedup over its serial baseline, compared to $0.24\times$ on CPU only. GPU strong
scaling efficiency reached $71.1\%$ at $P=2$ ranks without degrading at $P \geq 4$, while the weak scaling 
showed a $6\times$ improvement at $P=8$. These results using a consumer GPU represent a lower bound.

\textbf{3. Communication masking is validated.} The range of the metric $M = T_{\text{accel}} / T_{\text{comm}}$ 
was between 12--9,000, which confirms the classical computation has masked MPI overhead by 1--3 orders of magnitude.

\textbf{4. QPU integration as an architectural proof-of-concept is validated.} 
10 VQE iterations were successfully completed on IBM \texttt{ibm\_marrakesh}, validating transpilation, cloud submission, 
PUB batching, and MPI coordination. The error $0.97\;E_h$ is expected given the limit of 10 iterations (the availability within IBM's free access plan)
and the introduction of shot-noise, not to an architectural limitation within the stack, however, the observed energy trend was consistently descending. 
A final GPU accelerated QPU run confirmed this triple-heterogeneous path's validity.

\textbf{5. Robustness confirmed through stress testing.} Checkpoint resilience, QPU latency spikes 
(random 0.5--2.0\,s delays), and midrun checkpoint deletion recovery tests all passed without issues 
of data loss or synchronization failure occurring.

\subsection*{Limitations and Conclusion}
The primary limitations found were the $O(2^n)$ per rank statevector cost, limited bandwidth when using consumer GTX 1650 Mobile GPU
(128 GB/s), the HWE ansatz lacking particle number conservation
(issue was mitigated through near-zero initialization but only fully resolvable with UCCSD), and conducting single 
deterministic runs (seed=42) on one Docker host.

Despite the limitations, this work demonstrates that a middleware focused on orchestration, 
distribution, and latency masking produces measurable speedups even on commodity hardware,  
complements other existing circuit-level frameworks such as XACC. Validation on the lower bound of accessible hardware confirms the 
distributed hybrid VQE is feasible without data-center HPC resources for applications in the fields of pharmaceutical 
drug discovery and materials science.

\end{document}
